% Transcription — fonds Alexandre Grothendieck, shelfmark 156-3, batch 1 (pages 1-20).
% Established from the facsimile archives/batches/156-3/batch-01.pdf (a mirror of
% https://grothendieck.umontpellier.fr/156-3.pdf; no cover sheet in the batch, so
% PDF page k = folder page k), per the skill transcribe-grothendieck. The folder
% is forty pages; this is the first of two batches.
%
% All twenty leaves are mathematics in his hand and all are transcribed; no page
% is skipped and none is another's text to be summarised. His own pagination
% 1-20 runs across the top of the leaves and coincides with the archivists'.
% Page 1 carries the heading « Réseaux via découpages. (8 Juin 86) » and, in a
% triangle in the top-right corner, « GF III » (Géométrie des formes, chapter
% III, after folders 156-1 and 156-2). No other date on these twenty pages. The
% inventory's bracketed « [Chapitre] » is the archivists'; the title is his.
%
% What the batch is about. He wants a new axiomatics for a « réseau » (network)
% structure on a set L of « lieux », built from « découpages ». Pages 1-4: the
% data are (a) a set S_0 ⊂ P_2(L) of « 0-sphères » (pairs of lieux in moderate
% relative position, « disjoints »), whence the finite moderate parts Modf(L);
% (b) for each T ∈ Modf(L) an equivalence relation R_T on the moderate
% complement C(T), whose classes are the connected components; (c) for each
% component A a frontier Ȧ ⊂ T. He expects the axioms to make (c) follow from
% (b), and (b) follow from the R_ε for ε ∈ S_0 alone. Example 1: the discrete
% structure (S_0 = P_2(L), R_T discrete, Ȧ = ∅), with an axiom that a maximal
% finite moderate part is L. Pages 5-6: « provisional » finiteness axioms Dec_1
% (each U has finitely many components) and Dec_2 (for U ⊂ V cofinite moderate,
% each component of U lies in one of V, Ȧ ∩ T ⊂ Ḃ); Example 2, the « ligne
% ordonnée », where Modf(L) is the finite totally ordered parts and the
% components of C(T) are the n+1 open intervals. Pages 7-10: « tronçons » of
% the line and their bords; the definition of a « tronçon ordonné » (axioms a),
% a'), b) filtering, c) divisibility, d) card > 1), segments / intervals /
% half-intervals, ∂L, S_0(L) = Drap_2(L) (flags of length 2) and
% Modf(L) = Drap(L), the components of C(S). Pages 11-15: a lemma that the
% « binordonné » structure is recovered from the ternary relation R(a,b,c)
% « b strictly between a and c », via a four-point lemma (x ≤ u ≤ v ≤ y) proved
% on page 13; then that R(a,b,c) is itself recovered from Drap_2(L) and the
% component decompositions of the C(ε) — b is singled out in {a,b,c} by
% X_ε ∩ X_ε' = ∅ for the two pairs containing it (« On gagne ! »), using
% divisibility. Page 16: remarks — the largest full subcategory of ordered sets
% with the stability properties is that of tronçons; what conditions R(a,b,c)
% must satisfy is left open. Pages 17-20: closed sub-tronçons and their ∂I,
% moderate relative position, closed moderate parts (a theorem on page 18:
% uniqueness of the decomposition into sub-tronçons and Ẋ = ∪ Ṫ_i),
% intersections and the moderate union X ∪_mod Y with (X ∪_mod Y)˙ ⊂ Ẋ ∪ Ẏ, the
% example [a,b] ∪_mod [b,c] = [a,c], a criterion for total order, and « Partie
% loc. fermée modérée »: C_mod(Y,X) and the least closed moderate Z containing
% it, written X −_mod Y in the margin. The page-20 statement is heavily
% overwritten and only partly read; it presumably continues in batch 2.
%
% Reading hazards fixed once here. L is his script L (\mathcal{L}); his script
% S of S_0 is set \mathfrak{S}_0 (page 1 writes it plain, S_0); P_2, P_f are
% \mathfrak{P}. « Drap » is his abbreviation (drapeaux), « Modf » his, as is
% « binordonné » (read so on pages 11, 16, 18). His circled letters and numbers
% (a), (b), (c), (1), (2) are set in parentheses. His underlining is set in bold.
% He writes the boundary of a part with a dot (Ȧ, Ẋ), set \dot{}; on page 17
% it is ∂I. The drawn 0-sphere (two dots joined by a line) that stands for
% « ε ∈ S_0 » on pages 3 and 4 is described in a note. The hand is fast
% drafting throughout; slanted left-margin notes on pages 1, 2, 7, 9, 11, 17,
% 18, 19 were turned and read only in part.
%
% Readings flagged and not settled: most of the bracket on page 1 (« En termes
% de l'axiomatique par les segments… »); the first line of page 7; « trois »
% for four elements in the page-13 lemma; the closing statement of page 20; the
% brackets [s_1, s_2] on page 10 (open intervals expected); « s ≠ t » on
% page 2 where the page-1 margin has « a = b ». Variances on the page and not
% repaired: « S ⊂ P_f(L) » for S ∈ P_f(L) (p. 9); t ∈ L_{∉ε} and {ε, t} for
% t ∈ L ∖ ε and ε ∪ {t} (p. 14).
%
% Pass: Opus 5.5 (claude-opus-5-5), 2026-09-24 — first pass, unchecked against the pages by a human.

\documentclass[11pt,a4paper]{article}
% Le préambule commun à toutes les transcriptions du fonds Grothendieck.
%
% Il tient en une idée : **l'incertitude se compose, elle ne se résout pas.**
% Une transcription qui lisse un mot douteux a détruit la seule chose pour
% laquelle on la faisait — un lecteur doit pouvoir savoir, à chaque endroit,
% ce qui était sur le feuillet et ce qui a été ajouté. Les six macros qui
% suivent sont donc l'appareil critique, et non de la décoration.
%
% Les mêmes commandes sont reconnues par `scripts/render.mjs`, qui produit la
% vue de lecture du site. Une macro ajoutée ici doit l'être aussi là-bas,
% faute de quoi le rendu échouera bruyamment — ce qui est voulu : un rendu
% silencieusement faux est pire qu'un rendu absent.

\ProvidesPackage{grothendieck}[2026/08/08 Transcriptions du fonds Grothendieck]

\usepackage{amsmath,amssymb,amsthm}
\usepackage{mathrsfs}
\usepackage{stmaryrd}
% Les diagrammes commutatifs sont partout dans ce fonds : c'est la langue
% même de la géométrie algébrique de Grothendieck, et une transcription qui
% les rendrait en prose ne transcrirait rien.
\usepackage{tikz-cd}
% Français par défaut — c'est la langue des feuillets — mais l'anglais est
% chargé aussi : la traduction et le résumé sont des documents anglais, et la
% typographie française (espaces avant les deux-points) y serait une faute.
% Ces documents basculent par \selectlanguage{english} en tête de corps.
\usepackage[english,french]{babel}
\usepackage{fontspec}
\usepackage{geometry}
\geometry{a4paper,margin=25mm,marginparwidth=28mm}
\usepackage{marginnote}
\usepackage{xcolor}
% ulem, not soul. `soul`'s \st and \ul are built on a character-by-character
% scan that XeLaTeX's font machinery breaks: the document fails with an
% undefined control sequence rather than degrading. ulem does the same two jobs
% and is Unicode-engine safe. `normalem` stops it hijacking \emph, which the
% transcriptions use for Grothendieck's own emphasis.
\usepackage[normalem]{ulem}
\usepackage{hyperref}
\hypersetup{colorlinks=true,linkcolor=black,urlcolor=black,pdfborder={0 0 0}}

% --- Métadonnées du lot -----------------------------------------------------
% Reprises telles quelles par le script de rendu, qui les affiche en tête de
% la vue de lecture. Elles fixent ce que le fichier prétend couvrir : sans
% elles, une transcription flotte sans point d'attache dans un fonds de seize
% mille feuillets.

\newcommand{\folder}[1]{\gdef\grothFolder{#1}}
\newcommand{\batch}[1]{\gdef\grothBatch{#1}}
\newcommand{\leaves}[2]{\gdef\grothFirst{#1}\gdef\grothLast{#2}}
\newcommand{\foldertitle}[1]{\gdef\grothFoldertitle{#1}}
\gdef\grothFolder{?}\gdef\grothBatch{?}\gdef\grothFirst{?}\gdef\grothLast{?}\gdef\grothFoldertitle{}

% --- Le résumé --------------------------------------------------------------
%
% Il ouvre la lecture modernisée, sous le titre « Résumé », et il est composé
% en petit. Ce n'est pas un ornement : c'est le seul passage du document qui
% ne soit pas des mathématiques, et un lecteur doit pouvoir le reconnaître
% avant de l'avoir lu — puis le sauter s'il connaît déjà le sujet, ou s'y
% arrêter s'il ne le connaît pas. Le filet qui le suit marque la reprise du
% corps.

\newenvironment{resume}
  {\small\setlength{\parskip}{0.45em}}
  {\par\medskip\hrule\medskip}

% --- Mots-clés ---------------------------------------------------------------
%
% En anglais, en clôture du résumé de la lecture modernisée : le vocabulaire
% moderne sous lequel chercher ce que les feuillets construisent. Le manifeste
% du site (`npm run manifest`) les extrait pour étiqueter la cote entière —
% cette ligne est la source unique des tags, il n'y a pas de fichier à côté.
\newcommand{\keywords}[1]{\par\smallskip\noindent
  {\footnotesize\textsc{Keywords} --- \textit{#1}}\par}

% --- L'appareil critique ----------------------------------------------------

% Le numéro de feuillet, en marge. C'est l'ancre qui permet de régler un
% désaccord sur une formule en regardant *un* feuillet plutôt qu'une chemise
% de six cents. Le site s'en sert aussi pour tourner les pages du fac-similé
% quand on fait défiler la transcription.
\newcommand{\leaf}[1]{%
  \marginnote{\footnotesize\color{gray!70}\textsf{#1}}%
  \label{leaf:#1}\ignorespaces}

% Illisible : on ne devine pas. Un mot inventé qui se lit comme les autres est
% le pire résultat possible.
\newcommand{\ill}{\textcolor{red!60!black}{[\,\ldots\,]}}

% Lecture douteuse : le mot est donné, et signalé comme incertain.
\newcommand{\uncertain}[1]{\uline{#1}}

% Ajout de l'éditeur — une lettre manquante, un mot sous-entendu.
\newcommand{\add}[1]{\textcolor{blue!50!black}{[#1]}}

% Biffé par Grothendieck lui-même. Ce qu'il a rayé fait partie du document :
% une rature dit souvent où la pensée a changé de direction.
\newcommand{\struck}[1]{\sout{#1}}

% Note du transcripteur, sur l'état matériel du feuillet.
\newcommand{\note}[1]{\par\smallskip{\small\color{gray!60!black}\textit{[#1]}}\par\smallskip}

% Note marginale de Grothendieck, distincte du corps du texte.
\newcommand{\marginal}[1]{\par\smallskip\noindent\hspace{1em}%
  {\small\color{gray!60!black}#1}\par\smallskip}

% --- Titre ------------------------------------------------------------------

\AtBeginDocument{%
  \begin{center}
    {\large\bfseries Fonds Alexandre Grothendieck --- cote n\textsuperscript{o}~\grothFolder}\\[2pt]
    {\itshape\grothFoldertitle}\\[4pt]
    {\small feuillets \grothFirst--\grothLast\ (lot \grothBatch)}\\[2pt]
    {\footnotesize\color{gray!60!black}Transcription établie d'après le fac-similé numérisé
     par l'Université de Montpellier.\\
     Les lectures incertaines sont soulignées, les illisibles notées [\,\ldots\,],
     les ajouts entre crochets.}
  \end{center}
  \vspace{1em}}

\endinput


\folder{156-3}
\batch{1}
\pages{1}{20}
\dating{1986}
\watermark{Édition de démonstration}
\foldertitle{[Chapitre] III. Réseaux via découpages : notes manuscrites (08/06/1986)}

\begin{document}
\page{1}
\section{Réseaux via découpages}

\note{titre de sa main, souligné, en tête de la page 1, suivi de « (8 Juin 86) » ; dans l'angle supérieur droit, dans un triangle, « GF III »}

\marginal{\struck{\ill{}}}
\note{dans l'angle supérieur gauche, deux ou trois mots obliques biffés d'un trait épais, non lus}

Je voudrais définir une \uncertain{autre axiomatique} des \struck{réseaux,} \struck{\ill{}} structures de réseau sur un ens. $\mathcal{L}$ (ens. des « lieux »).

(a) $\mathfrak{S}_0$ \struck{\ill{}} $\subset \mathfrak{P}_2(\mathcal{L})$ l'ens. des \textbf{0-sphères} \struck{plongées} ou \uncertain{des} \struck{paires} \struck{\ill{}} (de lieux) modérées. [Si $\{a,b\} \in \mathfrak{S}_0$, on dit \struck{\ill{}} que $a$, $b$ sont \textbf{disjoints} ou en position relative modérée]\note{la phrase entre crochets est écrite au-dessus de la ligne}
\note{les lettres (a), (b), (c) sont cerclées sur la page, ici et plus loin ; de même (1), (2) à la page 16}
[En termes de l'axiomatique par les segments, cela signifie que \struck{(c'est-à-dire $a$, $b$ distincts, $a$, $b$ \ill{} ou bien)}, \ill{} que $\{a,b\} \in \mathfrak{P}_2(\mathcal{L})$, que ou bien $\{a,b\} \cap \mathcal{L}_S \neq \emptyset$ ($a$ ou $b$ est un \uncertain{nœud} du réseau) ou bien \add{dit} $a$, $b$ (si $a$, $b$ \struck{sont} \uncertain{lignes}) \ill{} : \uncertain{sur des} lignes différentes, sont, s'ils sont \uncertain{sur une même} ligne, \ill{} $\{a,b\} = \partial I$, pour $I \in S$].
\marginal{Mais on dit \uncertain{aussi} que $a$, $b$ \uncertain{sont en} position rel. modérée si $a = b$ \textbf{ou} $\{a,b\} \in \mathfrak{S}_0$}
\note{la note précédente est écrite en oblique dans la marge gauche. Le crochet « En termes de… » n'est lu qu'en partie : plusieurs mots sont surchargés}

(b) Une partie \textbf{finie} $T$ est dite « \textbf{modérée} » \struck{\ill{}} si \struck{$\forall a, b \in T$, \ill{}} $\mathfrak{P}_2(T) \subset \mathfrak{S}_0$ i.e. deux éléments quelc. de $T$ sont en position relative modérée (\uncertain{sont} 2 à 2 « disjoints »).

\page{2}
\struck{De plus, on veut que la structure implique}

\struck{(b) $\forall T \in P_{f\,\mathrm{mod}}(\mathcal{L})$, on se donne une relation d'équivalence $R_T$ dont \ill{}}
\note{ces lignes sont biffées de quatre traits obliques ; le (b) cerclé est hachuré}

Si $T$ est une partie finie modérée, on \struck{désigne par} pose
\[
  C_{\mathrm{mod}}(T) = \{t \in \mathcal{L} \smallsetminus T \mid \{T \cup \{t\}\} \text{ modérée}\}
\]
\ill{} des $t \in \mathcal{L} \smallsetminus T$ en \ill{} v. à v. de $T$.

[\ill{} partie finie modérée \ill{} en position relative \ill{} si leur réunion est \ill{}. $t$ est dit en pos. rel. \ill{} v. à v. de $T$ si $\{t\}$ l'est ; $s$ et $t$ sont dits en pos. rel. modérée si $\{s,t\}$ est modérée, i.e. \struck{ou} \uncertain{$s \neq t$} ou $\{s,t\} \in \mathfrak{S}_0$…]
\note{le début de ces lignes est recouvert par les notes de la marge. On attendrait « $s = t$ ou », comme dans la marge de la page 1}

\marginal{\textbf{le complémentaire modéré de $T$}}
\marginal{Question : $T$ \ill{} connaît \ill{} $C(T)$ \ill{} \textbf{NB} $C(T')$ \ill{} $C(T)$ ? Oui \ill{}}
\marginal{Quant à $C_{\mathcal{L}}(T)$ ? \ill{} Si $T \subset T'$ \ill{} l'inverse \ill{}. Donc, si $t \in T'$, $t \notin T$, \struck{\ill{}} \struck{($T \cup \{t\}$ modérée)} \ill{} $t \in C(T)$ \ill{}}
\note{trois notes obliques couvrent la marge gauche de la page, en partie biffées ; elles ne sont lues que par fragments}

(b) $\forall\, T \in \mathrm{Modf}(\mathcal{L})$\note{sous « Modf », de sa main : « parties mod. finies »}, une relation d'équivalence $R_T$ dans $C(T)$ (dont les classes s'appellent les \textbf{comp. connexes} de $C(T)$)
\marginal{\textbf{NB} \ill{} \ill{} $C(T)$}

\page{3}
(c) $\forall\, T \in \mathrm{Modf}(\mathcal{L})$, et $A \in C(T)/R_T$ une classe suivant $R_T$, une partie $\dot{A} \subset T$, appelée \textbf{frontière} de $A$.

Ces données (a) (b), (c) étant assujetties à satisfaire certains axiomes.

Je prévois que les axiomes impliquent \struck{Soit \ill{} $Z \subset C(T)$, \ill{} pour} ceci
\begin{itemize}
  \item[1°)] que la donnée (c) résulte déjà de (b)
  \item[2°)] que la donnée (b) résulte de la donnée des $R_T$, pour le cas particulier où $T \in \mathfrak{S}_0$.
\end{itemize}
Donc la seule donnée essentielle serait (a) celle des 0-sphères (i.e. la notion de \struck{pos.} lieux en pos. relative modérée, condition \add{réflexive et} symétrique)\note{« réflexive et » est écrit au-dessus de la ligne} et (b) la donnée, pour $\varepsilon \in \mathfrak{S}_0$, de la relation d'équivalence $R_\varepsilon$ dans $C(\varepsilon)$.
\note{devant « $\varepsilon \in \mathfrak{S}_0$ », un petit dessin : deux points joints par un trait, une 0-sphère}

\struck{Donc on}

\textbf{Axiomes provisoires :}

\struck{Soit $X$} Appelons \uncertain{les} $U$ de $\mathcal{L}$ \uncertain{formés} \ill{} $C(T)$ \textbf{parties modérées cofinies}. On

\page{4}
veut que si $U$ \struck{$(= C(T))$} est à la fois modérée finie et modérée cofinie, donc si $S$ et $T$ sont deux parties modérées finies \struck{telles} disjointes telles que \struck{tout} a)\note{« a) » est écrit au-dessus du mot biffé} $S \cup T$ modérée i.e. \ill{} $s \in S$ modéré \uncertain{par rapport à} $t \in T$, et b) Tout $u \in \mathcal{L}$ modéré \uncertain{par rapport à} $S$ est dans $T$, \struck{donc $S \cup T = \mathcal{L}$} et \add{finie}
\note{devant « $s \in S$ », le même dessin de 0-sphère qu'à la page 3}
i.e. $S \cup T$ est partie \uncertain{finie} modérée maximale, donc
\[
  S \cup T = \mathcal{L},
\]
et $\mathfrak{S}_0 = \mathfrak{P}_2(\mathcal{L})$, enfin les relations d'équivalence sont discrètes, et les $\dot{A}$ sont vides. \uncertain{Pour un} \ill{} \ill{},

\ill{} définit un réseau, \ill{}

\textbf{Ex 1} \struck{Déf} \textbf{structure discrète} sur $\mathcal{L}$
\note{« définit un réseau » est relié par un trait à « discrète »}
\[
  \mathfrak{S}_0 = \mathfrak{P}_2(\mathcal{L}) \Longrightarrow
  \left\{
  \begin{aligned}
    &\mathrm{Modf}(\mathcal{L}) = \mathfrak{P}_f(\mathcal{L}) \\
    &C_T = \mathcal{L} \smallsetminus T \quad \text{si } T \in \mathfrak{P}_f(\mathcal{L})
  \end{aligned}
  \right.
\]
\struck{$C(T) = C$}
\[
  R_T = \text{relation d'équivalence discrète dans } \mathcal{L} \smallsetminus T,
\]
donc les compos. $A$ sont réduites à un pt \struck{Pour Donc}
\[
  \dot{A} = \emptyset
\]

\textbf{Axiome} Si $T$ est partie finie modérée maximale, alors $T = \mathcal{L}$, $\mathcal{L}$ est discret.

\page{5}
On va faire des hypothèses de finitude draconiennes, peut-être provisoires
\marginal{Finitude}

\textbf{Dec$_1$} L'ens. des comp. connexes de tout $U$ est \textbf{fini}.

\struck{Définition Partie modérée fermée}

\textbf{Dec$_2$} Soient $U$ \add{$= C(S)$, $V = C(T)$}\note{écrit au-dessus de la ligne} parties modérées cofinies, avec $U \subset V$ i.e. $T \subset S$. Alors toute composante connexe $A$ \add{de $U$} est contenue dans une comp. connexe (unique bien sûr) $B$ de $V$. De plus, on a
\[
  \underbrace{\dot{A}}_{\subset\, S} \cap\, T \subset \underbrace{\dot{B}}_{\subset\, T}
\]

\page{6}
\textbf{Ex 2} Soit \struck{\ill{}} $\mathcal{L}$ un ens. ordonné, on suppose $\mathcal{L}$ filtrant croissant, filtrant décroissant, sans plus grand ni plus petit élément, loc. filtrant croissant et filtrant décroissant, divisible.

On appelle un tel ensemble une \textbf{ligne ordonnée}.
\[
  \mathfrak{S}_0 \subset \mathfrak{P}_2(\mathcal{L})
\]
\struck{défini par $\{x,y\} \in \mathfrak{P}_2(\mathcal{L})$} $=$ parties \struck{qui} de card. 2 qui sont \textbf{tot.} ordonnées
\[
  \mathrm{Modf}(\mathcal{L}) = \text{parties finies tot. ordonnées}
\]
\struck{i.e. telles que}

Soit $T \in \mathrm{Modf}(\mathcal{L})$, $T = \{t_1, t_2, \ldots, t_n\}$, $t_1 < t_2 < \cdots < t_n$.

Alors les composantes connexes sont
\[
  \underbrace{]-\infty, t_1[}_{t_1},\ \underbrace{]t_1, t_2[}_{t_1,\, t_2},\ \ldots,\ \underbrace{]t_{n-1}, t_n[}_{t_{n-1},\, t_n},\ \underbrace{]t_n, +\infty[}_{t_n}
\]
\marginal{Pts frontière}
($n+1$ composantes connexes)

\page{7}
Cette structure est donc \uncertain{déterminée} par sa \struck{\ill{}} \uncertain{donnée} binaire \uncertain{par} \ill{}, i.e. elle ne dépend que \uncertain{du} binaire.

\struck{Partie modérée}

\struck{\ill{}}

\textbf{Tronçon} de la ligne : partie de la forme
\[
  [s,t] \quad (s \leq t), \qquad ]-\infty, s] \text{ ou } [s, +\infty[, \qquad \mathcal{L} = ]-\infty, +\infty[
\]
\textbf{Bord} du tronçon : $\{s,t\}$ dans le premier cas, $\{s\}$ dans le second, $\emptyset$ \add{dans le} troisième.

\marginal{\ill{} ligne : \uncertain{tout tronçon} \ill{} \uncertain{fermé} \ill{} \uncertain{tronçon} \ill{}. \uncertain{Le binaire} \ill{} \uncertain{unité}. D'autre part, \uncertain{le bord d'un} \ill{} \uncertain{Bord} \uncertain{défini} \ill{}}
\note{note oblique dans la marge gauche, à côté des intervalles, lue en partie seulement}

\textbf{Deux tronçons} sont dits en position relative modérée, si leurs bords \ill{}. \textbf{Partie modérée} : réunion \add{finie}\note{« finie » est écrit au-dessus de la ligne} de tronçons \uncertain{mutuellement} modérés, et disjoints. \struck{Deux} Par définition, le bord est la réunion (disjointe) des bords.

\page{8}
\struck{Proposition}

\struck{\textbf{Définition} Segments, intervalles, demi-intervalles (cas d'une structure de réseau). \struck{Segment}}
\note{au-dessus de « segments, intervalles, demi-intervalles », une accolade porte « tronçons » ; la ligne entière est biffée de trois traits obliques}

\textbf{Définition} Un ens. ordonné $\mathcal{L}$ est appelé un \textbf{tronçon ordonné} si
\begin{itemize}
  \item[a)] $\forall a \in \mathcal{L}$ tel que $\mathcal{L}_{>a}$ \struck{$=$} $\neq \emptyset$, $\mathcal{L}_{>a}$ est filtrant décroissant \struck{i.e. $\forall x, y$ $a, y$ tels que} [si $a, y, x \in \mathcal{L}$ tels que] $x > a$, $y > a$, alors $\exists z \in \mathcal{L}$ avec $x, y \geq z > a$
  \item[a')] Axiome dual : si $a, x, y \in \mathcal{L}$ tels que $x < a$, $y < a$, alors $\exists z \in \mathcal{L}$ avec $x, y \leq z < a$
  \item[b)] $\mathcal{L}$ est filtrant croissant et filtrant décroissant i.e. $x, y \in \mathcal{L} \Longrightarrow \exists a, b \in \mathcal{L}$ avec $a \leq x, y \leq b$
  \item[c)] Si $x, y \in \mathcal{L}$ tels que $x < y$, $\exists z$ tel que $x < z < y$. \struck{\ill{}}
  \item[d)] \struck{\ill{}} $\operatorname{card} \mathcal{L} > 1$
\end{itemize}
\note{la partie entre crochets dans a) est écrite au-dessus de la ligne, en remplacement du passage biffé}
\marginal{\textbf{NB} les cond. a), b), c) \uncertain{satisfaites} si card $\leq 1$ !}

On dit que $\mathcal{L}$ est un \textbf{segment} si $\mathcal{L}$ admet un plus petit et un plus grand élément, un \textbf{intervalle} s'il n'admet ni l'un ni l'autre,

\page{9}
un \textbf{demi-intervalle} s'il admet un plus petit élément mais pas un plus grand, ou bien l'inverse.

\struck{L'oppo}

\textbf{NB} L'opposé d'un tronçon, segment, intervalle, demi-intervalle est idem. L'opposé d'un ½ intervalle fermé à g. est ½ fermé à droite, et inversement. \struck{Les} \textbf{Tout} tronçon est l'un des trois (segment, ½ intervalle, intervalle) à l'exclusion \ill{} deux à \uncertain{deux}.
\marginal{\ill{} \uncertain{tronçons} \ill{} \uncertain{binaires}}
\marginal{tr. \ill{} \textbf{binaires}}
\note{deux notes obliques dans la marge gauche, en face du NB ; la première n'est pas lue}

Si $\mathcal{L}$ est un tronçon, on pose \struck{$\partial\mathcal{L}$}
\[
  \partial \mathcal{L} = \{x \in \mathcal{L} \mid x \text{ est plus petit élément ou plus grand élément}\}
\]

Soit $\mathcal{L}$ un tronçon, et
\[
  \begin{aligned}
    \mathfrak{S}_0(\mathcal{L}) = \mathrm{Drap}_2(\mathcal{L}) &= \{\varepsilon \in \mathfrak{P}_2(\mathcal{L}) \mid \varepsilon \text{ totalement ordonné}\} \\
    &= \{\{x,y\} \mid x < y\}
  \end{aligned}
\]
\[
  \text{Déf. :=} \qquad \mathrm{Modf}(\mathcal{L}) = \mathrm{Drap}(\mathcal{L}) = \{S \subset \mathfrak{P}_f(\mathcal{L}) \mid S \text{ tot. ord.}\}
\]
\note{« Drap » : drapeaux. Le signe $\subset$ devant $\mathfrak{P}_f(\mathcal{L})$ est sur la page, là où l'on attendrait $\in$}

\page{10}
Si $S \in \mathrm{Modf}(\mathcal{L})$, on a donc \struck{$C(S) =$}
\[
  C(S) = \{t \in \mathcal{L} \mid t \notin S,\ S \cup \{t\} \text{ tot. ord.}\}
\]
On a dit comment on décomposait en « composantes connexes » un $C(S)$, $S = \{s_1, \ldots, s_n\}$, $s_1 < s_2 < \cdots < s_n$. Les comp. connexes sont les ens. non vides parmi
\[
  \mathcal{L}_{<s_1},\ [s_1, s_2],\ \ldots,\ [s_{n-1}, s_n],\ \mathcal{L}_{>s_n}
\]
\note{les crochets des intervalles intermédiaires sont lus tels qu'ils paraissent (le premier est surchargé) ; on attend des intervalles ouverts}

\begin{itemize}
  \item[] [cas $n = 0$ : une seule comp. $\mathcal{L}$
  \item[] cas $n = 1$ : \struck{une} deux composantes $\mathcal{L}_{<s}$, $\mathcal{L}_{>s}$ (celles qui sont non vides)
  \item[] cas $n \geq 2$ : voir ce \uncertain{qui est} plus \uncertain{haut}]
\end{itemize}

\ill{} non vides \uncertain{intervalles} étant les seules \ill{} $\mathcal{L}_{<s_1}$ (si $s_1$ plus petit élément) et $\mathcal{L}_{>s_n}$ (si $s_n$ plus grand élément). Les bords \struck{sont} respectivement
\[
  \underbrace{\{s_1\}}_{\text{si comp. existe}},\ \{s_1, s_2\},\ \ldots,\ \{s_{n-1}, s_n\},\ \underbrace{\{s_n\}}_{\text{si comp. existe}}
\]
\note{après le premier $\{s_1\}$, une surcharge biffée}

Je dis que l'on récupère la structure \struck{\ill{}} binaire du tronçon (donc aussi la structure (a) (b) (c)), en termes de la seule structure (a), + (b) pour $T \in \mathfrak{S}_0$.

\page{11}
Pour ceci, on note \struck{premier \ill{}} d'abord\note{« d'abord » est écrit au-dessus de la ligne}

\textbf{Lemme} Une structure de \struck{binaire} tronçon binordonné est connue quand on connaît la relation sur \struck{$\{a,b,c\}$} trois éléments $a, b, c$
\[
  \mathcal{R}(a,b,c) \qquad \text{$b$ est \textbf{strictement entre} $a$ et $c$}
\]
(au sens que $a < b < c$ \textbf{ou} $a > b > c$)
\note{un triple trait vertical, dans la marge gauche du corps, met en évidence la ligne de $\mathcal{R}(a,b,c)$}
\marginal{relation \ill{} \uncertain{symétrique} en $a$, $c$}
\marginal{\struck{On la connaît de} \textbf{NB} Cette relation $\mathcal{R}$ équivaut à la relation $\overline{\mathcal{R}}$ non stricte \struck{$a < b < c$ ou} $a \leq b \leq c$ ou $a \geq b \geq c$ \struck{$a \neq b \neq c$ ; $\{a,b,c\}$} $b$ entre $a$ et $c$ ($\Longleftrightarrow$) \ill{} ou $\{a,b,c\}$ est un \ill{} cardinal $\leq 2$}
\note{notes obliques, en partie biffées, dans la marge gauche}

\textbf{Dém} \struck{Choisissons} On en déduit la connaissance de\note{ces mots sont écrits au-dessus de la ligne} $\mathfrak{S}_0(\mathcal{L})$ (i.e. de $\mathrm{Drap}_2(\mathcal{L})$) \ill{} \uncertain{formulation} $a$ et $b$ « disjoints » i.e. distincts et comparables i.e. $\{a,b\} \in \mathrm{Drap}_2(\mathcal{L})$ pour le binaire\,:

$a$ et $b$ disj\supplied{oin}ts ssi $\exists\, c$ tel que $c$ \uncertain{soit strictement entre} $a$ et $b$ (utiliser divisibilité \struck{\ill{}} \uncertain{du} binaire)

\struck{\ill{}} \textbf{Choisissons} $\varepsilon \in \mathfrak{S}_0$ (existe, car $\exists\, x, y \in \mathcal{L}$, $x \neq y$, \struck{si $x, y \in$} par hyp., \ill{} $a \leq x, y \leq b$, et on a $a \leq b$ et $a \neq b$ \struck{car} \ill{} on aurait $x = y$)
\marginal{\ill{} \uncertain{en utilisant} \ill{} \uncertain{arbitraire} \ill{}}

\page{12}
et choisissons une « origine » $a$ dans $\varepsilon = \{a, b\}$ (ce qui va déterminer la relation d'ordre). \struck{Soit} \struck{Alors} Considérons sur $\mathcal{L}$ \uncertain{la relation d'ordre} \ill{} définissant \uncertain{le binaire}, telle que $a < b$. On va voir qu'elle est entièrement déterminée par la \uncertain{conn.} du \struck{\ill{}} couple\note{« couple » est écrit au-dessus de « paire », biffé} $(a,b)$.

On connaît
\[
  \left.
  \begin{aligned}
    \mathcal{L}_{\geq b} &= \{x \in \mathcal{L} \mid x = b \text{ ou } b \text{ entre } a \text{ et } x\} \\
    \mathcal{L}_{\leq a} &= \{x \in \mathcal{L} \mid x = a \text{ ou } a \text{ entre } b \text{ et } x\} \\
    ]a, b[ &= \{x \in \mathcal{L} \mid x \text{ entre } a \text{ et } b\}
  \end{aligned}
  \right\}
\]
\marginal{déterminés par $\mathcal{R}$}
\note{« $x = b$ ou », « $x = a$ ou » sont écrits au-dessus de la ligne ; au-dessus du premier, un « $x \geq b$ ou » biffé}

donc \struck{Quand} Soient $u, v \in \mathcal{L}$, quand a-t-on $u \leq v$ ? Il faut et il suffit qu'il existe \struck{$x \in$} $x, y \in \mathcal{L}$ tels que
\[
  x \in \mathcal{L}_{\leq a},\quad y \in \mathcal{L}_{\geq b} \quad \text{et} \quad x \leq u \leq v \leq y
\]
i.e.
\begin{itemize}
  \item[] \struck{$\{x, y, u, v\}$ est $\in \mathrm{Drap}(\mathcal{L})$} $x \in \mathcal{L}_{\leq a}$, $y \in \mathcal{L}_{\geq b}$ (donc $x < y$)
  \item[] \struck{$x, y \in$} \struck{avec $x$ plus petit élément}
  \item[] \struck{$u$ est entre} $u$, $v$ sont entre $x$ et $y$
  \item[] $u$ est entre $x$ et $v$
  \item[] $v$ est entre $u$ et $y$
\end{itemize}
\note{les trois dernières conditions sont réunies par un crochet à gauche}

\page{13}
Donc il faut simplement vérifier le

\textbf{Lemme} Soit $\mathcal{L}$ un ens. ordonné, et $x, y, u, v$ \uncertain{trois}\note{sic ; quatre éléments sont nommés} éléments, tels que
\begin{itemize}
  \item[a)] $u$, $v$ sont entre $x$ et $y$
  \item[b)] $u$ est entre $x$ et $v$
  \item[c)] [$v$ est entre $u$ et $y$ \struck{\ill{}}]\note{sur la page, le crochet ouvrant précède « c) »}
  \item[d)] $x < y$,
\end{itemize}
\note{b) et c) sont réunis par une accolade : « une des deux suffit »}
alors on a
\[
  x \leq u \leq v \leq y
\]

\textbf{Dém} \struck{\ill{}} des conditions b), c), d) \ill{} pour $I = \{x, y, u, v\}$\note{une surcharge biffée entre $I$ et le signe $=$} est un drapeau, \struck{Prenons} et OPS $\mathcal{L} = I$ est un ens. tot. ordonné \struck{par $x <$ \ill{}} à 2, 3, ou 4 éléments. \struck{Alors Prenons \ill{}}
\struck{Comme $x$ est le plus petit élément, donc \ill{} (\ill{} $x < y$), pour l'un \ill{} \ill{} suppose $u < x \,(< y)$, ou $u < x$, mais $v < x$ \ill{} comme $u$ est entre $x$ et $v$, on aurait \ill{} et $x < y$, $v < u < x < y$}
\note{ce passage est encadré et biffé de plusieurs traits obliques ; il n'est lu qu'en partie}
\ill{} $x \leq u, v \leq y$, \struck{il reste} on doit avoir \struck{à prouver $u \leq v$} et comme $u$ est entre $x$ et $v$, on doit avoir $x \leq u \leq v$. OK

\textbf{NB} On a utilisé a) la divisibilité de la rel. d'ordre, b) le fait que $\mathcal{L}$ soit filtrant croissant et filtrant décroissant.

Et a) n'a servi qu'à montrer que la relation « $a$ et $b$ disjoints » est déductible de « $c$ est str. entre $a$ et $b$ ».

\page{14}
Si on se donne à la fois les deux relations
\[
  \left\{
  \begin{aligned}
    &\mathcal{R}(a,b) &&\{a,b\} \in \mathrm{Drap}_2(\mathcal{L}) \\
    &\mathcal{R}(a,b,c) &&b \text{ est str. entre } a \text{ et } c
  \end{aligned}
  \right.
\]
alors pour récupérer le binaire, il suffit de supposer $\mathcal{L}$ filtrant croissant et filtrant décroissant \struck{et \ill{} filtrant cr. et décroiss. \ill{} \ill{} compatible \ill{}}
\note{deux lignes et demie biffées de traits horizontaux et obliques}

\struck{Dém} On va voir que la connaissance de\note{écrit au-dessus de la ligne : « que la connaissance de $\mathrm{Drap}_2(\mathcal{L})$, et \uncertain{puis} »} $\mathrm{Drap}_2(\mathcal{L})$, et \uncertain{puis} de la décomp. de $\varepsilon \in \mathrm{Drap}_2(\mathcal{L})$\note{le $\mathcal{L}$ de cette parenthèse est surchargé}, de la décomp. de
\[
  C(\varepsilon) = \{t \in \mathcal{L}_{\notin \varepsilon} \mid \{\varepsilon, t\} \in \mathrm{Drap}_3(\mathcal{L})\}
\]
en composantes connexes, implique la connaissance de la relation
\[
  \mathcal{R}(a,b,c) \ : \ b \text{ strict entre } a \text{ et } c.
\]
\note{l'indice de $\mathcal{L}$ et l'accolade $\{\varepsilon, t\}$ sont lus tels qu'ils paraissent ; on attend $t \in \mathcal{L} \smallsetminus \varepsilon$ et $\varepsilon \cup \{t\}$}

Première condition : $\{a,b,c\} \in \mathrm{Drap}_3(\mathcal{L})$ (connue de la donnée de $\mathrm{Drap}_2(\mathcal{L})$)

\note{ici, un croquis : un segment portant les points $a$, $b$, $c$ dans cet ordre}

Si $b$ est entre $a$ et $c$, regardons les
\[
  \mathcal{J}_{abc} = \text{composante de } C(\{a,b\}),
\]

\page{15}
trois composantes (relatives à la donnée de $\{a,b,c\}$)
\[
  \begin{aligned}
    X_{\{a,b\}} &= (\text{comp. de } C(\{a,b\}) \text{ contenant } c) = \mathcal{L}_{>b} \\
    X_{\{b,c\}} &= (\text{comp. de } C(\{b,c\}) \text{ contenant } a) = \mathcal{L}_{<b} \\
    X_{\{c,a\}} &= (\text{comp. de } C(\{c,a\}) \text{ contenant } b) = ]a, c[
  \end{aligned}
\]
\note{sur la page, entre chaque $X$ et le premier signe $=$, une première notation est biffée et illisible ; les deuxième et troisième lignes répètent « comp. de » et « contenant » par des traits}

On trouve
\[
  \left.
  \begin{aligned}
    X_{\{a,b\}} \cap X_{\{b,c\}} &= \emptyset \\
    X_{\{a,b\}} \cap X_{\{c,a\}} &= ]b, c[ \neq \emptyset \\
    X_{\{b,c\}} \cap X_{\{c,a\}} &= ]a, b[ \neq \emptyset
  \end{aligned}
  \right\} \text{divisibilité !}
\]

donc $b$ est distingué, dans $J = \{a, b, c\}$, par la propriété que pour les deux \struck{parties} paires $\varepsilon, \varepsilon' \in \mathfrak{P}_2(J)$ qui le contiennent, on a
\[
  X_{\varepsilon} \cap X_{\varepsilon'} = \emptyset.
\]
On gagne !

\textbf{NB} On a à nouveau utilisé la divisibilité.

\page{16}
\textbf{Remarques} (1) \struck{Considérons} Sur la notion de tronçon. On choisit une sous-cat. pleine $\mathcal{T}$ de celle des ens. ordonnés, satisfaisant
\begin{itemize}
  \item[a)] Si $\mathcal{L} \in \operatorname{Ob} \mathcal{T}$, $\mathcal{L}$ est filtrant croissant et filtrant décroissant, et de cardinal $\neq 1$ \struck{non vide} (donc $\geq 2$)
  \item[b)] Si $x, y \in \mathcal{L} \in \operatorname{Ob} \mathcal{T}$, et $x < y$, $\exists\, z \in \mathcal{L}$, $x < z < y$ (divisibilité)
  \item[c)] Si $x \in \mathcal{L} \in \operatorname{Ob} \mathcal{T}$, \struck{donné} et si $\mathcal{L}_{>x} \neq \emptyset$, alors $\mathcal{L}_{>x} \in \operatorname{Ob} \mathcal{T}$.
  \item[c')] \textbf{Dual} de c)
\end{itemize}
Il y a un plus grand $\mathcal{T}$ qui satisfait ces conditions, et $\mathcal{T}$ est stable par passage \uncertain{à l'ordre} opposé. Les objets de $\mathcal{T}$ sont les tronçons ordonnés.

(2) Il faudrait expliciter quelles conditions la relation $\mathcal{R}(a,b,c)$ doit satisfaire, pour définir une structure de \struck{tr.} tronçon binordonné, et itou pour la donnée de $\mathfrak{S}_0 \subset \mathfrak{P}_2(\mathcal{L})$ et les $\mathcal{R}_\varepsilon$ ($\varepsilon \in \mathfrak{S}_0$) sur les $C(\varepsilon)$…

\page{17}
\textbf{Sous-tronçon} \add{fermé, $I$}\note{« fermé, $I$ » est écrit au-dessus de la ligne} d'un tronçon $\mathcal{L}$ : sous-ens. de la forme $[a,b]$ (avec $a < b$) ou $\mathcal{L}_{\leq a}$ (avec $a$ pas le plus petit élément) ou $\mathcal{L}_{\geq b}$ (avec $b$ pas un plus grand élément), \add{ou $\mathcal{L}$}\note{écrit au-dessus de la ligne} : ce sont donc des tronçons, pour la structure induite. On pose $\partial I = \{a,b\}$ resp. $\{a\}$ resp. $\{b\}$, resp. $\emptyset$. \struck{Deux sous Deux} Deux sous-tronçons fermés sont dits en pos. rel. modérée si leurs frontières le sont, i.e. si \ill{} frontière de l'un \ill{} \uncertain{sous-tronçon fermé} \ill{} : soit \ill{} de l'autre.
\marginal{\ill{}}
\note{plusieurs notes obliques dans la marge gauche, non lues}

\struck{\ill{} Deux \ill{} tronçons \ill{} ssi leurs frontières \ill{} \uncertain{fermés}, \uncertain{larges} \ill{} \uncertain{en pos. gén.} \ill{} \uncertain{modérée} \ill{} \uncertain{intersection} \ill{}. Alors leur intersection \ill{} tronçon \ill{} dont la frontière \ill{} des frontières.}
\note{trois lignes et demie biffées de longs traits, lues par fragments}

\ill{} tronçon \uncertain{fermé} : \ill{} partie de $\mathcal{L}$ \ill{}

\textbf{Partie fermée modérée} de $I$ :

[toute partie qui est de la forme de réunion d'une famille finie de sous-tronçons fermés mutuellement disjoints et mutuell. \ill{} \uncertain{en pos. rel. modérée}]

\uncertain{Prenons} \ill{} \ill{} \uncertain{ayant comme base} les $]a, b[$ pour $a < b$, et \uncertain{les} $\mathcal{L}_{>a}$ et $\mathcal{L}_{<a}$. Alors \struck{\ill{}} pour \uncertain{tout} sous-\uncertain{tronçon} fermé $I$, $\partial I$ est la frontière.
\marginal{\textbf{NB} \ill{}}

\page{18}
\textbf{Th} Soit $\mathcal{L}$ un tronçon binordonné, soit $X \subset \mathcal{L}$, tel qu'il existe $I \subset \mathrm{Pstrf}(\mathcal{L})$ (ens. des ps-tronçons fermés de $\mathcal{L}$), $I$ fini, \struck{\ill{}} tel que
\begin{itemize}
  \item[a)] $\forall\, i, j \in I$, $i \neq j$,
  \[
    \left\{
    \begin{aligned}
      &T_i \cap T_j = \emptyset \\
      &T_i \text{ et } T_j \text{ en pos. rel. mod.}
    \end{aligned}
    \right.
  \]
  i.e. \struck{\ill{}} $\dot{T}_i \cup \dot{T}_j$ est partie finie modérée \uncertain{comme} $\forall\, a \in \dot{T}_i$, $b \in \dot{T}_j$, $a \neq b$, \ill{} $\{a,b\} \in \mathfrak{S}_0(\mathcal{L}) = \mathrm{Drap}_2(\mathcal{L})$\note{la lettre lue $\mathfrak{S}_0$ ressemble ici à un $B_0$}
  \item[] \struck{\ill{}}
  \item[b)] $X = \bigcup T_i$ ;
\end{itemize}
\note{il note ici le bord d'une partie par un point en exposant, $\dot{T}_i$ ; à la page précédente il écrivait $\partial I$}

Alors
\begin{itemize}
  \item[1°)] $I$ est unique [les $T_i$ sont les comp. conn. de $X$]
  \item[2°)] \struck{\ill{}} $\dot{X} = \bigcup_{i \in I} \dot{T}_i$

  (donc $\dot{X} = \emptyset$ ssi \struck{\ill{}} $X = \mathcal{L}$ ou $X = \emptyset$)
\end{itemize}
\marginal{\textbf{NB} $X = \emptyset$ ssi $I = \emptyset$}

On dit que $X$ est une \textbf{partie modérée fermée}.
\marginal{\ill{}}
\note{dans la marge gauche, une note verticale, non lue, qui commence peut-être par « \uncertain{Ex.} »}

\textbf{Déf} Si $X$, $Y$ sont deux parties fermées modérées, on dit que $X$ et $Y$ sont en pos. relative modérée si $\dot{X} \cup \dot{Y}$ \ill{}, i.e. $\forall\, x \in \dot{X}$, $y \in \dot{Y}$, $x \neq y$, on a $\{x,y\} \in \mathrm{Drap}_2$.

\page{19}
\textbf{Prop} \struck{\ill{}} Soient $X$, $Y$ deux parties fermées modérées, en pos. rel. modérée. Alors $X \cap Y$ est une partie \add{fermée}\note{écrit au-dessus de la ligne} modérée, et
\[
  \boxed{(X \cap Y)^{\cdot} \subset \dot{X} \cup \dot{Y}}
\]
\struck{De} $X$ et $Y$ ont une borne sup. \struck{dans l'ensemble des parties} \ill{} des parties modérées \struck{qui contiennent $X$} \struck{\ill{} $X$ et $Y$}. On note cette partie $X \cup_{\mathrm{mod}} Y$, et on a
\[
  \boxed{(X \cup_{\mathrm{mod}} Y)^{\cdot} \subset \dot{X} \cup \dot{Y}}
\]
\marginal{\ill{} \uncertain{Dém} \ill{} \uncertain{convenons} \ill{} \uncertain{modérée} \ill{}}
\marginal{\textbf{NB} \ill{} une partie \ill{} \uncertain{plus petits} \ill{} \uncertain{modérés} \ill{} Ex. \uncertain{$A = \{0, \tfrac12, \tfrac13, \ldots\}$} \ill{}}
\note{les deux marges portent de longues notes obliques, lues par fragments seulement}

Soit $(X_i)_{i \in I}$ une famille finie de parties fermées modérées, mutuellement en pos. rel. \uncertain{modérée}. Alors $\bigcap X_i$ est une partie modérée, et les $X_i$ ont une borne sup. dans l'ens. des parties fermées modérées. On a
\[
  \Bigl(\bigcup_{\mathrm{mod},\, i} X_i\Bigr)^{\cdot} \subset \bigcup_i \dot{X}_i
\]

\textbf{Ex}
\[
  [a,b] \cup_{\mathrm{mod}} [b,c] = [a,c]
\]
mais en général
\[
  [a,b] \cup [b,c] \neq [a,c]
\]
(prenant $a$ plus petit élément, $c$ plus grand élément, \ill{} \struck{\ill{}} $Z$, $b$ \struck{\ill{}} \uncertain{non} comparable à tout élément de $\mathcal{L}$. Si c'est

\page{20}
vrai pour tout $b$, alors $\mathcal{L}$ est tot. ordonné. Plus généralement, pour que $\mathcal{L}$ soit tot. ordonné, il faut et il suffit que pour tout $b \in \mathcal{L}$, $b$ soit comparable à tout élément, ce qui s'écrit
\[
  \underbrace{\mathcal{L}_{\leq b} \cup_{\mathrm{mod}} \mathcal{L}_{\geq b}}_{\mathcal{L}} = \mathcal{L}_{\leq b} \cup \mathcal{L}_{\geq b}
\]

\subsection*{Partie loc. fermée modérée}

\textbf{Th} Soient $X$ \struck{\ill{}}, $Y$ deux parties, $X \supset Y$ (deux fermées modérées, \struck{avec} $X$ et $Y$ en pos. relative modérée)
\struck{Alors Considérons \ill{} une plus petite partie $Z$ fermée modérée contenant $X \smallsetminus Y$.}
\struck{Ex} \textbf{Déf}\note{« Déf » est écrit en surcharge sur « Ex »}
\[
  C_{\mathrm{mod}}(Y, X) = \{x \in X \mid x \notin Y,\ x \text{ en pos. rel. mod. avec } Y\}
\]
Alors $\exists$ plus petit ens. fermé modéré $Z$ \struck{$\supset A$} \uncertain{$\subset X$, et} \ill{} $C_{\mathrm{mod}}(Y, X)$, \ill{} partie fermée modérée $T \subset X$, $T \cap Y = \emptyset$, $T$ en pos. rel. \uncertain{mod. avec} $Y$, \struck{De plus $Z$} \uncertain{$Z$ est en position} \uncertain{rel.} modérée \ill{} $Y$.
\marginal{On \uncertain{pose} $Z = X \mathbin{-_{\mathrm{mod}}} Y$}
\note{la fin de la page est surchargée d'ajouts interlinéaires et n'est lue qu'en partie}

\end{document}
